\documentclass[11pt,a4paper]{article}
\usepackage{isabelle,isabellesym}

\usepackage[T1]{fontenc}
\usepackage[utf8]{inputenc}

\usepackage{amsmath}
\usepackage{amssymb}

% this should be the last package used
\usepackage{pdfsetup}

\urlstyle{rm}
\isabellestyle{literal}

% Provide any Isabelle symbols not defined by isabellesym.sty.
\providecommand{\isasymto}{\ensuremath{\to}}
\providecommand{\isasymrightarrow}{\ensuremath{\rightarrow}}
\providecommand{\isasymsubseteq}{\ensuremath{\subseteq}}
\providecommand{\isasymmp}{\ensuremath{\mp}}
\providecommand{\isasympm}{\ensuremath{\pm}}
\providecommand{\isasymne}{\ensuremath{\neq}}

% Session / entry names, defined in sessionnames.tex.
% The AFP entry / session names, defined once here and used throughout the
% document, so that a rename is one edit.
\newcommand{\sessMTTM}{\texttt{Multitape\_TM\_Substrate}}
\newcommand{\sessAE}{\texttt{Multitape\_Alphabet\_Enlargement}}
\newcommand{\sessAR}{\texttt{Multitape\_Alphabet\_Reduction}}
\newcommand{\sessART}{\texttt{Multitape\_Alphabet\_Roundtrip}}
\newcommand{\sessNTR}{\texttt{Nondeterministic\_Tape\_Reduction}}
\newcommand{\sessDTR}{\texttt{Deterministic\_Tape\_Reduction}}
\newcommand{\sessNTH}{\texttt{Nondeterministic\_Time\_Hierarchy}}


\begin{document}

% \sloppy: raise the line-breaking tolerance so LaTeX rearranges breaks
% around the long unbreakable \isa{} identifier tokens (in text blocks and
% proof displays) rather than letting them overrun the right page edge.
% Verified to clear every Overfull \hbox >= 20pt in this document.
\sloppy

\title{A Nondeterministic Multitape Turing Machine Substrate\\[6pt]
       \large Formal Foundations in Isabelle/HOL}
\author{Andr\'as Z.~Salamon \and Michael Wehar}
\maketitle

\begin{abstract}
A self-contained, reusable formalisation of nondeterministic
multitape Turing machines in Isabelle/HOL, depending only on
\texttt{HOL-Library}.

A machine's transition is a \emph{relation}, not a function: it can
allow several moves at each step, so a machine may follow many different
computations on one input --- the machine is \emph{nondeterministic}.
The machine accepts when at least one computation reaches an accepting
state.  Deterministic machines, which have a single move everywhere, are
the special case.  The number of
tapes is stored as ordinary data rather than fixed in the machine's
type, so a single theorem can speak about machines with any number of
tapes.  Hierarchy theorems need exactly this.

The entry also justifies its least conventional modelling choice.  Each
tape has a marker at its left end, and the model lets a machine move
that marker to the right, discarding a stretch of tape it can never
return to.  We prove that this changes nothing the machine computes: the
run before the move reaches an accepting state exactly when the matching
run after it does.  This turns an assumption that
such models usually leave implicit into a proved fact, and it gives a
reusable tool for arguments that relocate or trim a tape.

On top of the model, the entry provides the constructions that
developments built on it share: facts about runs that each of them would
otherwise re-prove, transport of a machine along injections on its states
and alphabet, widening and exchanging tapes, running a machine as one
factor of a composite, restoring the endmarker discipline, an interface
for exact unary clocks, and running a machine alongside such a clock.

This entry is the shared foundation for several companion developments:
alphabet enlargement (\sessAE) and alphabet reduction (\sessAR), tape
reduction, and the time hierarchy theorems.  It is also self-contained
and usable on its own.

The model began as an adaptation of Dalvit and Thiemann's multitape
Turing machine formalisation~\cite{Dalvit2022:verified} (the AFP entry
\texttt{Multitape\_To\_Singletape\_TM}), and extends it in two ways: the
number of tapes is part of a machine's data rather than fixed in its
type; and a machine is itself an ordinary value that theorems can range
over, rather than a fixed context --- a locale --- inside which each
theorem is proved.  Like their formalisation,
a machine's transition is a relation, so machines may be
nondeterministic; that sets both developments apart from the Archive's
other Turing-machine entries, whose transitions are functions.  The
names differ from theirs (\texttt{ktm} for \texttt{mttm},
\texttt{ktc\_config} for \texttt{mt\_config}, \texttt{move} for
\texttt{dir}, and so on), so that a theory can import both entries.
\end{abstract}

\paragraph*{Acknowledgement.}
This formalisation was developed with proof engineering by
Anthropic's Claude Opus~4.5, 4.6, 4.7, and 4.8 as the work progressed,
and revised with Opus~5.5, together with other Claude family models in
supporting roles.
Claude Code was used to structure the work and for access to Isabelle/HOL.
We also thank Arthur Freitas Ramos, who provided feedback, contributed
several tooling improvements, and collected some constructions over
the model into a shared session, while working with OpenAI Codex
driven by OpenAI's GPT-6 and by Claude Opus~4.8.


\tableofcontents

\newpage

\section{Overview}

\subsection{The substrate}

This session defines a reusable model of nondeterministic multitape
Turing machines, self-contained and usable on its own.  A machine's
transition is a \emph{relation}, not a function: a machine may have
several possible moves at each step, so nondeterminism is the default,
and a deterministic machine is the special case in which every situation
offers at most one move.  Dalvit and Thiemann model the transition
relationally in the same way; two further design choices set this
substrate apart from their \texttt{Multitape\_To\_Singletape\_TM}, and
are fixed here once.

First, the number of tapes is part of a machine's data rather than
fixed in its type.  A single theorem can then range over machines with
any number of tapes.

Second, a machine is itself an ordinary value (a piece of data), so
it can be built, compared, listed one after another, and even encoded as
a word and fed to another machine as input.  A hierarchy theorem needs
exactly this last ability: it builds one machine that runs every other
machine in turn and disagrees with each.  A model that instead fixes a
single machine as a background setting, and proves each theorem inside
that setting, can still state facts that hold for every machine, but it
cannot handle a machine as data in this way.

The model and its predicates are defined in
\texttt{Multitape\_Substrate\_Core.thy}, and \texttt{Multitape\_Substrate.thy}
proves their basic metatheory.  Companion theories add the origin-float justification (\texttt{Multitape\_Origin\_Float.thy},
described below) and the finite-control and time-bound convention layers
the speedup entries reuse (\texttt{Multitape\_Finite\_Control.thy},
\texttt{Multitape\_Finite\_Patch.thy}, \texttt{Multitape\_Time\_Convention.thy}).  The companion
\sessAE{} and \sessAR{} entries reuse this vocabulary unchanged, and their
overviews take it as given.

Two other Turing-machine developments in the Archive of Formal Proofs
sit alongside this substrate.  Both give a machine a single move at each
step.  Balbach's \texttt{Cook\_Levin}~\cite{Balbach2023:Cook} formalises
multitape machines, with the number of tapes a value as here, for the
Cook--Levin theorem and $\mathrm{NP}$-completeness.  Xu, Zhang, Urban,
and Joosten's \texttt{Universal\_Turing\_Machine}~\cite{Xu2019:universal}
builds a single-tape universal machine for computability, without any
timing analysis.  In neither may a machine take several moves at once,
so nondeterminism has to be added on top rather than being built in.
That is the main way this substrate differs from both.

The constructions over the model are described in the last subsection of this overview.

\subsection{Machines, configurations, and runs}

\paragraph{Machines.}  A machine is a value of the datatype
$\mathrm{ktm}$, the ten-component descriptor
$\mathrm{KTM}\ Q\ \Sigma\ \Gamma\ bl\ le\ \delta\ s\ t\ r\ k$: a state
set $Q$, an input alphabet $\Sigma$ and a tape alphabet $\Gamma$ (both
ordinary \emph{value} sets over an element type ${'a}$), a blank $bl$
and left-end marker $le$ (two special \emph{tape} symbols that validity
keeps in $\Gamma$ but not $\Sigma$, so both are usable on a tape yet
never as input), a transition relation $\delta$, start,
accept
and reject states $s, t, r$, and a tape count $k :: \mathbb{N}$.  The
components are read by the projections \texttt{Q\_ktm},
\texttt{Sigma\_ktm}, \texttt{\isasymGamma\_ktm}, \texttt{bl\_ktm},
\texttt{le\_ktm}, \texttt{delta\_ktm}, \texttt{s\_ktm}, \texttt{t\_ktm},
\texttt{r\_ktm}, \texttt{k\_ktm}.

\paragraph{Configurations.}  A configuration is a value of
$\mathrm{ktc\_config}$, written $\mathrm{Config}_{\mathrm{M}}\ q\ ts\ n$:
the current state $q$, the tape contents $ts$ (a function
$\mathbb{N} \to \mathbb{N} \to {'a}$ giving the symbol in cell $j$ of
tape $i$ as $ts\,i\,j$), and the head positions $n$ (a function
$\mathbb{N} \to \mathbb{N}$).  Tapes and cells are indexed by
$\mathbb{N}$; in a valid run, every tape beyond the machine's count $k$
stays blank.  The start configuration is
$\mathrm{init\_config\_ktm}\, M\, w$: state $s$, tape $0$ holding the
endmarker $le$ at cell $0$ followed by the input word $w$, every other
in-range tape holding just $le$, and all heads at cell $0$.

\paragraph{The step relation.}  $\mathrm{ktm\_step}\, \delta$ is the
one-step relation on configurations, parametrised by $\delta$ alone.
From $\mathrm{Config}_{\mathrm{M}}\, q\, ts\, n$, let the
\emph{read vector} $\lambda k.\ ts\,k\,(n\,k)$ collect the symbol under
each head; if
$\bigl(q,\ \lambda k.\ ts\,k\,(n\,k),\ q',\ a,\ move\bigr) \in \delta$,
the machine may step to $\mathrm{Config}_{\mathrm{M}}\, q'\, ts'\, n'$,
where $ts'$ overwrites cell $n\,k$ of each tape $k$ with $a\,k$ and
$n'$ moves each head by $move\,k \in \{L, R, N\}$.  Because $\delta$ is
a relation, several tuples may match one configuration.  That is exactly
where nondeterminism enters.

\paragraph{Validity, language, and time.}  The remaining predicates,
the ones the transformation theorems carry as hypotheses and
conclusions, are built on top:
\begin{description}
\item[$\mathrm{valid\_ktm}\, M$] \emph{``$M$ is a legal machine''.}
  The structural invariant: $Q$ and $\Gamma$ finite,
  $\Sigma \subseteq \Gamma$, the states $s, t, r$ in $Q$, blank and
  endmarker tape-but-not-input symbols, $t \neq r$, $k > 0$, $\delta$
  correctly typed, the \emph{left-endmarker read discipline} (a
  transition reading $le$ rewrites $le$ and never moves left of it),
  and every tape beyond $k$ frozen blank.  Notably it does \emph{not}
  constrain how $le$ is \emph{written}; that is the separate
  $\mathrm{le\_unique}$ below, so a machine that plants a fresh $le$
  stays valid.
\item[$\mathrm{well\_formed\_ktm}\, M$] \emph{a legal machine that also
  starts cleanly and keeps its endmarker.}  $\mathrm{valid\_ktm}\, M$
  plus non-degeneracy ($s \neq t$, $s \neq r$, $le \neq bl$) and the
  endmarker write discipline $\mathrm{le\_unique}\, M$.  It is the
  hypothesis through which the alphabet-transformation theorems learn
  that the machine they simulate never writes a fresh endmarker.
\item[$\mathrm{le\_unique}\, M$] \emph{the left endmarker is never
  freshly planted.}  A machine property in the same family as
  $\mathrm{det\_ktm}$: no transition writes $le$ where it was not
  already reading it ($a'\, j = le \Rightarrow a\, j = le$ over
  $\delta$).  Folded into $\mathrm{well\_formed\_ktm}$.  The machine
  that achieves the $(1+\varepsilon)n$ speed-up in \sessAE{} is the
  one construction that forgoes it: it writes a fresh $le$ past the
  input it has consumed (see origin floating, below), so it is
  $\mathrm{valid\_ktm}$ but not $\mathrm{well\_formed\_ktm}$.
\item[$\mathrm{det\_ktm}\, M$] \emph{$M$ is deterministic.}  $\delta$
  is a partial function: each pair of a state and a read vector has at
  most one successor triple $(q', a, move)$.
\item[$L(M) = \mathrm{Lang\_ktm}\, M$] \emph{the language $M$ accepts.}
  The input words $w$ (with $\mathrm{set}\, w \subseteq \Sigma_M$) from
  which some run reaches an accepting configuration:
  \[
    (\mathrm{init\_config\_ktm}\, M\, w,\ c)
      \in (\mathrm{ktm\_step}\, \delta_M)^{*}
    \quad\text{with}\quad \mathrm{ktc\_state}\, c = t .
  \]
  The reflexive-transitive closure makes acceptance \emph{existential}
  over runs.
\item[$\mathrm{accepts\_in\_time\_ktm}\, M\, w\, m$] \emph{$M$ accepts
  $w$ within $m$ steps.}  The time-bounded variant: some run of length
  at most $m$ (an iterated step $(\mathrm{ktm\_step}\, \delta_M)^{\,n}$
  with $n \leq m$) reaches the accept state.
\end{description}
These six are exactly the entries the \sessAE{} and \sessAR{}
glossaries list under ``shared substrate predicates''; their
definitions live here.

\subsection{Scope and modelling choices}

This entry is \emph{foundational infrastructure}.  It fixes the machine
model and its invariants and proves no complexity results about specific
languages; those live in the companion entries.  Besides the constructions
described below, two of the facts its companions rely on are worth singling
out.  The first is the structural left-end-marker fact
(origin floating, below).  The second is the shared time-bound convention
vocabulary the linear-speedup entry consumes: the $\max(n + 2, \cdot)$
``convention'' predicate and its ``eventual'' companion, the
finite-control \texttt{time\_cleanup} that upgrades an eventual bound to a
convention bound on all inputs, and a minimal counterexample
\texttt{clean\_convention\_unattainable} showing that the clean
constant-free convention form is \emph{not} attainable in general --- a
valid machine can accept its language yet overshoot the convention floor
on its shortest input.  That is why the effective linear-speedup
statements in \sessAE{} keep an explicit additive constant or hold only
from a threshold on, and why its constant-free form for all inputs
decides the short inputs from a finite table.  Three modelling choices
are worth knowing before building on it.
\begin{itemize}
\item \emph{No separate input tape.}  The input word occupies tape
  $0$, an ordinary work tape, so \texttt{valid\_ktm} requires
  $\Sigma \subseteq \Gamma$ (every input symbol is also a tape symbol)
  and a transformation may re-encode the input \emph{in place} rather
  than copy it to a fresh tape.
\item \emph{One left-end marker, two rules.}  A single left-end marker
  $le$ obeys a \emph{read} rule: reading $le$ rewrites it and never
  moves off the left edge.  This rule is part of \texttt{valid\_ktm}.
  The stronger \emph{write} rule $\mathrm{le\_unique}$ ($le$ is never
  freshly planted) belongs instead to \texttt{well\_formed\_ktm}, so a
  machine that plants a fresh $le$ is still valid.  The
  $(1+\varepsilon)n$ speed-up machine of \sessAE{} is exactly such a
  machine.
\item \emph{One-way-infinite tapes.}  Cells are indexed by $\mathbb{N}$
  with the endmarker pinning cell $0$; every tape beyond the machine's
  count $k$ is frozen blank.
\end{itemize}

\subsection{Origin floating: the left-endmarker write, validated}

Of the choices above, moving the left-end marker is the one with real
proof content, and \texttt{Multitape\_Origin\_Float.thy} supplies it.  The model lets
a machine place a fresh marker at some cell $d$, which cuts off cells
$0$ to $d-1$.  That prefix is then unreachable, because the read rule
forbids the machine from ever stepping left of the marker.  Let
$\mathrm{shift\_rel}\, d$ be the operation that shifts every tape $d$
cells to the right, leaving the state and heads unchanged and moving the
symbol at cell $p$ to cell $p + d$.  The theory then shows that a run and
its shifted copy match step for step, in both directions and at every
level:
\begin{itemize}
\item single steps --- $\mathrm{ktm\_step\_shift\_forward}$ and
  $\mathrm{ktm\_step\_shift\_reverse}$;
\item whole runs --- $\mathrm{ktm\_relpow\_shift\_forward}$ and
  $\mathrm{ktm\_relpow\_shift\_reverse}$;
\item acceptance --- $\mathrm{shift\_reach\_state\_forward}$ and
  $\mathrm{shift\_reach\_state\_reverse}$: a run reaches a nominated state
  iff its shift does.
\end{itemize}
So placing a fresh marker neither gains nor loses behaviour.  Allowing
the machine to write this marker is therefore sound, and an assumption
that such a model would usually leave implicit becomes a theorem here.
The $(1+\varepsilon)n$ speed-up machine in \sessAE{} is the first
user of these lemmas: it moves the marker to reclaim the space taken up
by the consumed input.  The lemmas themselves are stated only in terms of
a single step of the machine, independently of that use, so they serve as
a reusable building block for any tape argument that relocates or trims a
tape, such as simulating a two-dimensional tape on ordinary tapes.

\subsection{Constructions over the model}

Every development built on the substrate composes machines, and each
composition changes the machine's type, its tape count, or its state
space.  The later theories collect the constructions that recur, so that
a development states its own construction and inherits the bookkeeping.
None of them builds a machine for a particular language.
\begin{description}
\item[\texttt{Multitape\_Run\_Facts}] Facts about the step relation
  and runs that mention no construction: acceptance and acceptance within
  a time bound (\texttt{Lang\_ktm\_iff}) with their introduction and
  elimination rules, the step relation under a
  larger or a union of transition relations, the halting states as dead
  ends, determinism lifted from transition tuples to runs, freezing a
  tape, and bounds on head movement.
\item[\texttt{Multitape\_Relabelling}] Transport of a machine along
  an injection on states and an injection on the tape alphabet
  (\texttt{relabel\_ktm}), which preserves the language up to the
  alphabet map (\texttt{relabel\_lang}) and the time within which a word
  is accepted (\texttt{relabel\_accepts\_in\_time} and its converse).
  This is what lets machines of arbitrary types be collected in one
  concrete type.  Two combinators on the tapes go with it: raising the
  tape count (\texttt{widen\_tapes}), which preserves the language
  (\texttt{widen\_lang}) and acceptance within a given time
  (\texttt{widen\_accepts\_in\_time}), and exchanging two tapes
  (\texttt{swap\_tapes}), whose runs are the original machine's runs
  with the two tapes exchanged (\texttt{swap\_run} and
  \texttt{swap\_run\_rev}).
\item[\texttt{Multitape\_Embedding}] Running a machine as one factor
  of a composite: a window of the composite's tapes and a region of its
  states (\texttt{embed\_delta}, \texttt{tag\_delta}), and the
  correspondences that read a factor's run as a run of the composite.
\item[\texttt{Multitape\_Endmarker\_Normal}] Restoring the endmarker
  write discipline $\mathrm{le\_unique}$ for any valid machine
  (\texttt{endmarker\_normal}).  The construction lifts the tape alphabet
  along \texttt{Some}, so the language is the original one mapped through
  \texttt{Some} (\texttt{endmarker\_normal\_lang}), and a word is accepted
  within the same time as before
  (\texttt{endmarker\_normal\_accepts\_in\_time} and its
  converse).
\item[\texttt{Multitape\_Clock\_Interface}] Exact unary clocks: a
  function $t$ is time-constructible when some well-formed deterministic
  machine reaches its accepting state from the unary input of length $n$
  in exactly $t(n)$ steps, with the facts about halting states and alarms
  that a clocked construction needs.
\item[\texttt{Multitape\_Clock\_Product}] A machine run alongside an
  alarm clock (\texttt{clock\_product}): the two machines sit on disjoint
  tapes, the composite writes the unary input for the clock and then
  alternates their steps, and it stops as soon as either one halts, so
  the clock bounds every run, accepting or not.  The theory proves the
  composite valid, with the endmarker write discipline, and relates its
  runs to the two factors' runs.
\item[\texttt{Multitape\_One\_Tape\_Clocks}] Concrete deterministic unary
  clocks shared by both hierarchy developments. The walk takes exactly
  $n+2$ steps; the marking sweep takes exactly $2n^2+3n+2$ steps, including
  two steps on empty input. Their transition systems and exact run proofs
  supply witnesses for \texttt{time\_constructible}.
\item[\texttt{Multitape\_Binary\_Codec}] A machine-free codec for
  writing a natural number as a fixed-width block of binary digits, with
  the two digit symbols as parameters (\texttt{nat\_to\_syms},
  \texttt{syms\_to\_nat}), its round trip and injectivity on the
  bounded range (\texttt{syms\_to\_nat\_nat\_to\_syms},
  \texttt{nat\_to\_syms\_inj}), and the width that indexes a finite set,
  $\lceil \log_2 |S| \rceil$ and at least one (\texttt{block\_width}).
\item[\texttt{Multitape\_Enum\_Type}] Machine-free support for
  applying a theorem quantified over enumerable types, such as the linear
  speedup theorem of \sessAE: a local copy of $\{0,\ldots,K-1\}$ is an
  instance of \texttt{enum}, so such a statement can be discharged with
  the \texttt{HOL-Types\_To\_Sets} local typedef (applied by the theory
  that uses it,
  so this entry does not depend on that session).
\end{description}

\newpage

% Isabelle's auto-generated session.tex contains the correct
% \input{} sequence in declaration order from this session's ROOT
% `theories` block.  A per-session PDF is self-scoped: it documents
% only Multitape_TM_Substrate's own theories (no document_theories inclusion).
\input{session}

\bibliographystyle{abbrv}
\bibliography{root}

\end{document}
